\section{Gaussian endpoint sufficiency and score gains}
\label{app:gaussian-extension}
\paragraph{Proof of Proposition~\ref{prop:ext-sufficiency}.}
\begin{proof}
Expanding the Gaussian density makes its parameter-dependent factor
$\exp\{s_Y A U-A^2J_{\rm H}/2\}$. Integrating over $A$ gives the first
posterior formula and shows that $U$ is sufficient for $Y$.
Likewise $S\mid Y,A\sim N(s_Y AM,C)$, giving the second formula.

The identity
\[
 \Delta=
 \left(\boldsymbol a-\frac MC\Sigma\boldsymbol1\right)^\mathsf T
 \Sigma^{-1}
 \left(\boldsymbol a-\frac MC\Sigma\boldsymbol1\right)
\]
proves nonnegativity and its equality condition. If $\Delta=0$, then
$U=(M/C)S$, so the posteriors agree. This also covers
$\boldsymbol a=0$, when both equal $\pi_0$.

For the converse, suppose $\Delta>0$ and define
$R=U-(M/C)S$. Conditional on $(Y,A)$, $R$ and $S$ are independent and
\begin{equation}
 R\mid Y,A\sim N(s_Y A\Delta,\Delta).
 \label{eq:ext-residual}
\end{equation}
Indeed, $\operatorname{Cov}(U,S\mid Y,A)=M$ and
$\operatorname{Var}(U\mid Y,A)=J_{\rm H}$. Thus, conditional on every finite
$S=s$, the distribution of $U$ has a strictly positive density on $\R$.
For $j>0$, differentiation under the integral is justified by Gaussian
damping in $\alpha$, and
\[
 \frac{d}{du}\log\frac{G_j(u)}{G_j(-u)}
 =\frac{G'_j(u)}{G_j(u)}+\frac{G'_j(-u)}{G_j(-u)}>0.
\]
Strictness follows from $\nu((0,\infty))>0$. Consequently $q_{\rm H}$ is not
constant conditional on $S$, and cannot agree almost surely with $q_{\rm E}$.
The proper-score differences are
\begin{align*}
 \E[\mathcal L_{\log}(q_{\rm E},Y)-\mathcal L_{\log}(q_{\rm H},Y)]
 &=\E\KL\!\left(\operatorname{Bern}(q_{\rm H})\,
       \middle\|\,\operatorname{Bern}(q_{\rm E})\right),\\
 \E[(Y-q_{\rm E})^2-(Y-q_{\rm H})^2]
 &=\E[(q_{\rm H}-q_{\rm E})^2].
\end{align*}
Both are therefore strictly positive. Since fixed $p_0$ makes $p_T$ an
invertible transformation of $S$, the same conclusion holds for the terminal
price.
\end{proof}
\paragraph{Conditional information--MMSE identity.}
To verify \eqref{eq:ext-conditional-immse}, use the Markov relation
$Y\mathbin{-}X\mathbin{-}W_\gamma$ to write
$I(Y;W_\gamma)=I(X;W_\gamma)-I(X;W_\gamma\mid Y)$.
Apply the information--MMSE identity to both terms. Their derivative is
\[
 \frac12\left\{\E\operatorname{Var}(X\mid W_\gamma)
       -\E\operatorname{Var}(X\mid W_\gamma,Y)\right\}
 =\frac12\E\operatorname{Var}\!\left(\E[X\mid W_\gamma,Y]\mid W_\gamma\right).
\]
The last conditional variance has the binary expression in
\eqref{eq:ext-conditional-immse}; it is bounded in expectation by
$\E X^2=\E A^2$. Finite second moments justify the integrated identity.

